\documentclass[crop,tikz]{standalone}

% Plots vector field and its magnitude of a finite solenoid of length
% h using the expressions from Martín-Luna et.al., "On the Magnetic
% Field of a Finite Solenoid", IEEE TRANSACTIONS ON MAGNETICS,
% VOL. 59, NO. 4, APRIL 2023

\usepackage{pgfplots}
\usetikzlibrary{arrows.meta}
\usepgfplotslibrary{colormaps}

\pgfplotsset{
  compat=1.18,
  inverted/.style = {
    every axis legend/.append style={
      draw=white,
      fill=hardblack,
      text=white
    }
  },
}
 
\begin{document}
\begin{tikzpicture}
  % Note that the expressions are written in terms of the
  % dimensionless ratio r/r_0. The x-axis shows the ratio x/r_0 and
  % the y-axis shows the ration z/r_0. I.e. the borders of the
  % solenoid are at x=-1 and x=1.
  \pgfmathsetmacro{\height}{2}; % ratio h/r_0
  \pgfmathsetmacro{\xmin}{-3};
  \pgfmathsetmacro{\xmax}{3};
  \pgfmathsetmacro{\ymin}{-3};
  \pgfmathsetmacro{\ymax}{3};
  \pgfmathsetmacro{\BO}{4*pi};
  \begin{axis}[
    xmin = {\xmin}, xmax = {\xmax},
    ymin = {\ymin}, ymax = {\ymax},
    domain = {\xmin}:{\xmax},
    domain y = {\ymin}:{\ymax},
    axis equal image,
    view = {0}{90},
    hide axis,
    colormap={whiteblue}{rgb=(1,1,1) rgb=(0,0,1)},
    declare function={
      % integrands (u in degree)
      Kint(\u,\kk) = 1/sqrt(max(1e-15, 1 - (\kk)^2*sin(\u)^2));
      Eint(\u,\kk) = sqrt(max(1e-15, 1 - (\kk)^2*sin(\u)^2));
      Piint(\u,\nn,\kk) = 1/(max(1e-15, 1 - (\nn)*sin(\u)^2)*sqrt(max(1e-15, 1 - (\kk)^2*sin(\u)^2)));
      % symmetric GL-Paar on [0,90] degree (mid=45, half=45)
      GLpairK(\xi,\wi,\kk) = (\wi)*(Kint(45 + 45*(\xi),\kk) + Kint(45 - 45*(\xi),\kk));
      GLpairE(\xi,\wi,\kk) = (\wi)*(Eint(45 + 45*(\xi),\kk) + Eint(45 - 45*(\xi),\kk));
      GLpairPi(\xi,\wi,\nn,\kk) = (\wi)*(Piint(45 + 45*(\xi),\nn,\kk) + Piint(45 - 45*(\xi),\nn,\kk));
      % complete elliptic integral K(k)
      ellipK(\kk) = (pi/4)*(
        GLpairK(0.9894009349916499, 0.027152459411754095, \kk) +
        GLpairK(0.9445750230732326, 0.062253523938647893, \kk) +
        GLpairK(0.8656312023878317, 0.095158511682492785, \kk) +
        GLpairK(0.7554044083550030, 0.12462897125553387,  \kk) +
        GLpairK(0.6178762444026437, 0.14959598881657673,  \kk) +
        GLpairK(0.4580167776572274, 0.16915651939500254,  \kk) +
        GLpairK(0.2816035507792589, 0.18260341504492359,  \kk) +
        GLpairK(0.09501250983763744,0.18945061045506850,  \kk)
      );
      % complete elliptic integral E(k)
      ellipE(\kk) = (pi/4)*(
        GLpairE(0.9894009349916499, 0.027152459411754095, \kk) +
        GLpairE(0.9445750230732326, 0.062253523938647893, \kk) +
        GLpairE(0.8656312023878317, 0.095158511682492785, \kk) +
        GLpairE(0.7554044083550030, 0.12462897125553387,  \kk) +
        GLpairE(0.6178762444026437, 0.14959598881657673,  \kk) +
        GLpairE(0.4580167776572274, 0.16915651939500254,  \kk) +
        GLpairE(0.2816035507792589, 0.18260341504492359,  \kk) +
        GLpairE(0.09501250983763744,0.18945061045506850,  \kk)
      );
      % complete elliptic integral Pi(n,k)
      ellipPi(\nn,\kk) = (pi/4)*(
        GLpairPi(0.9894009349916499, 0.027152459411754095, \nn, \kk) +
        GLpairPi(0.9445750230732326, 0.062253523938647893, \nn, \kk) +
        GLpairPi(0.8656312023878317, 0.095158511682492785, \nn, \kk) +
        GLpairPi(0.7554044083550030, 0.12462897125553387,  \nn, \kk) +
        GLpairPi(0.6178762444026437, 0.14959598881657673,  \nn, \kk) +
        GLpairPi(0.4580167776572274, 0.16915651939500254,  \nn, \kk) +
        GLpairPi(0.2816035507792589, 0.18260341504492359,  \nn, \kk) +
        GLpairPi(0.09501250983763744,0.18945061045506850,  \nn, \kk)
      );
      kk(\r,\z,\s) = sqrt(4*\r/((1 + \r)^2 + (\z + 0.5*\s*\height)^2)); % k_\pm
      kp(\r,\z) = kk(\r, \z, 1); % k_+
      km(\r,\z) = kk(\r, \z, -1); % k_-
      sigma(\r) = 4*\r/(1 + \r)^2;
      %
      Ipm(\r,\z,\k) = -4/sqrt(\r)*(ellipE(\k)/\k - (1 - 0.5*\k^2)*ellipK(\k)/\k); % I_\pm
      Jpm(\r,\z,\k,\s) = \k/sqrt(\r)*(\z + 0.5*\s*\height)*((1 - \r)/(1 + \r)*ellipPi(sigma(\r),\k) + ellipK(\k)); % I'_\pm
      Ip(\r,\z) = Ipm(\r,\z,kp(\r,\z)); % I_+
      Im(\r,\z) = Ipm(\r,\z,km(\r,\z)); % I_-
      Jp(\r,\z) = Jpm(\r,\z,kp(\r,\z),1); % I'_+
      Jm(\r,\z) = Jpm(\r,\z,km(\r,\z),-1); % I'_-
      %
      Bz0(\r,\z) = 0.5*\BO*((\z + 0.5*\height)/sqrt(1 + (\z + 0.5*\height)^2) - (\z - 0.5*\height)/sqrt(1 + (\z - 0.5*\height)^2)); % B_z(0,z)
      Br(\r,\z) = \r == 0 ? 0 : \BO/(4*pi)*(Im(\r,\z) - Ip(\r,\z)); % B_r(r,z)
      Brz(\r,\z) = \r == 0 ? Bz0(\r,\z) : \BO/(4*pi)*(Jp(\r,\z) - Jm(\r,\z)); % B_z(r,z)
      rr(\x,\y,\z) = sqrt(\x^2 + \y^2); % r(x,y,z)
      Bx(\x,\y,\z) = rr(\x,\y,\z) == 0 ? 0 : Br(rr(\x,\y,\z),\z)*cos(atan2(\y,\x)); % B_x(x,y,z)
      By(\x,\y,\z) = rr(\x,\y,\z) == 0 ? 0 : Br(rr(\x,\y,\z),\z)*sin(atan2(\y,\x)); % B_y(x,y,z)
      Bz(\x,\y,\z) = Brz(rr(\x,\y,\z),\z); % B_z(x,y,z)
      Ba(\x,\y,\z) = sqrt(Br(rr(\x,\y,\z),\z)^2 + Bz(\x,\y,\z)^2); % |\vec{B}(x,y,z)|
    },
    ]
    % magnitude of magnetic field
    \addplot3 gnuplot [
      raw gnuplot,
      surf,
      shader=interp,
      no marks,
    ] {
      set samples 101;
      set isosamples 101;
      set xrange [\xmin:\xmax];
      set yrange [\ymin:\ymax];
      set urange [\xmin:\xmax];
      set vrange [\ymin:\ymax];
      B0 = 4*pi;
      h = \height;
      k(r,z,s) = sqrt(4.*r/((1. + r)**2 + (z + 0.5*s*h)**2));
      kp(r,z) = k(r,z,1.);
      km(r,z) = k(r,z,-1.);
      sigma(r) = 4.*r/(1. + r)**2;
      Ipm(r,z,k) = -4./sqrt(r)*(EllipticE(k)/k - (1. - 0.5*k**2)*EllipticK(k)/k);
      Jpm(r,z,k,s) = k/sqrt(r)*(z + 0.5*s*h)*((1. - r)/(1. + r)*EllipticPi(sigma(r),k) + EllipticK(k));
      Ip(r,z) = Ipm(r,z,kp(r,z));
      Im(r,z) = Ipm(r,z,km(r,z));
      Jp(r,z) = Jpm(r,z,kp(r,z),1);
      Jm(r,z) = Jpm(r,z,km(r,z),-1);
      Bz0(r,z) = 0.5*B0*((z + 0.5*h)/sqrt(1 + (z + 0.5*h)**2) - (z - 0.5*h)/sqrt(1 + (z - 0.5*h)**2));
      Br(r,z) = r == 0 ? 0 : B0/(4*pi)*(Im(r,z) - Ip(r,z));
      Brz(r,z) = r == 0 ? Bz0(r,z) : B0/(4*pi)*(Jp(r,z) - Jm(r,z));
      r(x,y,z) = sqrt(x**2 + y**2);
      Bx(x,y,z) = r(x,y,z) == 0 ? 0 : Br(r(x,y,z),z)*cos(atan2(y,x));
      By(x,y,z) = r(x,y,z) == 0 ? 0 : Br(r(x,y,z),z)*sin(atan2(y,x));
      Bz(x,y,z) = Brz(r(x,y,z),z);
      Ba(x,y,z) = sqrt(Br(r(x,y,z),z)**2 + Bz(x,y,z)**2);
      splot "++" using 1:2:(Ba($1,0,$2))
    };
    % magnetic vector field
    \addplot3[
      red,
      -{Latex[scale=0.3]},
      samples=51,
      samples y=51,
      quiver={
        u={Bx(x,0,y)},
        v={Bz(x,0,y)},
        scale arrows=0.04,
      },
    ] {0};
    % solenoid
    \draw[very thick] (-1,{-\height/2}) -- (-1,{\height/2});
    \draw[very thick] (1,{-\height/2}) -- (1,{\height/2});
  \end{axis}
\end{tikzpicture}
\end{document}
